\subtituloA{Importancia de la multidisciplinariedad}

\lista{
  *Diseño gráfico:**** Conocimientos en diseño visual y experiencia de usuario para crear interfaces atractivas y funcionales. |
  \textbf{Programación:} Habilidades en lenguajes de programación tanto del lado cliente como del servidor. |
  \textbf{Gestión de proyectos:} Organización y planificación del desarrollo web en equipos multidisciplinarios. |
  \textbf{Marketing y SEO:} Optimización para motores de búsqueda y estrategias de posicionamiento web.
}

\nota{Para el examen}{
¿Cuál es el objetivo principal del W3C en el desarrollo web? La

\textbf{Respuesta:} es: estandarizar las tecnologías web para garantizar la interoperabilidad entre diferentes navegadores y dispositivos, asegurando que los documentos HTML, CSS y otros lenguajes se comporten de manera consistente en toda la web. La opción que afirma que el W3C solo crea navegadores es incorrecta porque su función es de estandarización, no de desarrollo de software. La opción que menciona que el W3C solo gestiona servidores es también falsa, ya que su alcance abarca todos los aspectos de las tecnologías web.

\vspace{2.5mm}\hrulefill\vspace{2.5mm}

¿Qué característica define a la web según el material? La

\textbf{Respuesta:} es: estar basada en el hipertexto, donde elementos conectados mediante referencias permiten acceder a información relacionada. La opción que describe la web como un sistema de correo electrónico es incorrecta. La opción que menciona que la web solo funciona con imágenes es falsa, ya que el hipertexto es su característica esencial. La opción que habla de una red de computadoras privada no representa la naturaleza abierta y pública de la web.
}